\documentclass[11pt,a4paper]{article}

% ====================================================================
%  Riepilogo consegna -- Gestione Utenti (sola visualizzazione)
%  Esame "Progettazione Full Stack" (Prof. Bianchini, UniBS)
%  Appello 20260629
%  Versione LaTeX del documento di studio.
%
%  Compilazione consigliata: pdflatex (2 passate per indice/riferimenti).
% ====================================================================

% ----------------------------- Pacchetti ----------------------------
\usepackage[utf8]{inputenc}
\usepackage[T1]{fontenc}
\usepackage[italian]{babel}
\usepackage[a4paper,margin=2.4cm]{geometry}
\usepackage{lmodern}
\usepackage{microtype}
\usepackage{xcolor}
\usepackage{listings}
\usepackage{booktabs}
\usepackage{tabularx}
\usepackage{longtable}
\usepackage{array}
\usepackage{enumitem}
\usepackage{seqsplit}
\usepackage{fancyhdr}
\usepackage{titlesec}
\usepackage{parskip}
\usepackage[hidelinks,colorlinks=true,linkcolor=black!70,urlcolor=blue!60!black]{hyperref}

% ----------------------------- Colori -------------------------------
\definecolor{codebg}{RGB}{248,248,248}
\definecolor{codeframe}{RGB}{220,220,220}
\definecolor{codekw}{RGB}{0,90,160}
\definecolor{codestr}{RGB}{160,60,30}
\definecolor{codecom}{RGB}{100,120,100}
\definecolor{codenum}{RGB}{120,120,120}
\definecolor{accent}{RGB}{40,70,120}
\definecolor{okgreen}{RGB}{30,130,60}
\definecolor{kored}{RGB}{180,40,40}

% ------------------------- Stili dei titoli -------------------------
\titleformat{\section}{\Large\bfseries\color{accent}}{\thesection}{0.6em}{}
\titleformat{\subsection}{\large\bfseries\color{accent!85!black}}{\thesubsection}{0.6em}{}
\titleformat{\subsubsection}{\normalsize\bfseries}{\thesubsubsection}{0.6em}{}

% ------------------------- Intestazioni -----------------------------
\pagestyle{fancy}
\fancyhf{}
\lhead{\small Gestione Utenti --- sola visualizzazione}
\rhead{\small Appello 20260629}
\cfoot{\thepage}
\renewcommand{\headrulewidth}{0.4pt}

% --------------------- Listings: stile generico ---------------------
\lstdefinelanguage{TypeScriptX}{
  morekeywords={
    import, from, export, const, let, var, function, async, await,
    return, if, else, for, while, interface, type, enum, class,
    extends, implements, public, private, readonly, new, void,
    Promise, string, number, boolean, null, undefined, true, false,
    useState, useEffect, useNavigate, this, default, throw
  },
  sensitive=true,
  morecomment=[l]{//},
  morecomment=[s]{/*}{*/},
  morestring=[b]',
  morestring=[b]",
  morestring=[b]`
}

\lstset{
  basicstyle=\ttfamily\footnotesize,
  backgroundcolor=\color{codebg},
  frame=single,
  rulecolor=\color{codeframe},
  keywordstyle=\color{codekw}\bfseries,
  stringstyle=\color{codestr},
  commentstyle=\color{codecom}\itshape,
  numberstyle=\tiny\color{codenum},
  numbers=left,
  numbersep=8pt,
  showstringspaces=false,
  breaklines=true,
  breakatwhitespace=false,
  columns=fullflexible,
  keepspaces=true,
  tabsize=2,
  extendedchars=true,
  literate=%
    {à}{{\`a}}1 {è}{{\`e}}1 {é}{{\'e}}1 {ì}{{\`i}}1 {ò}{{\`o}}1 {ù}{{\`u}}1
    {À}{{\`A}}1 {È}{{\`E}}1 {É}{{\'E}}1 {ç}{{\c c}}1
    {👤}{{[user]}}1 {🚧}{{[WIP]}}1 {✅}{{Si}}1
}

% Stile "diagramma" per i blocchi ASCII (WebML)
\lstdefinestyle{diagram}{
  basicstyle=\ttfamily\scriptsize,
  backgroundcolor=\color{codebg},
  frame=single,
  rulecolor=\color{codeframe},
  numbers=none,
  keywordstyle=\relax,
  commentstyle=\relax,
  stringstyle=\relax,
  breaklines=false,
  columns=fullflexible,
  keepspaces=true,
  extendedchars=true,
  literate=%
    {à}{{\`a}}1 {è}{{\`e}}1 {é}{{\'e}}1 {ì}{{\`i}}1 {ò}{{\`o}}1 {ù}{{\`u}}1
}

% Comandi abbreviati
% \file: percorsi di file -- spezzabili a fine riga (seqsplit) per non sforare il margine.
\newcommand{\file}[1]{\texttt{\small\seqsplit{#1}}}
\newcommand{\code}[1]{\texttt{\small #1}}
\newcommand{\si}{\textcolor{okgreen}{\textbf{S\`i}}}

% ====================================================================
\begin{document}

% --------------------------- Frontespizio ---------------------------
\begin{titlepage}
\centering
\vspace*{2cm}
{\Huge\bfseries\color{accent} Gestione Utenti\par}
\vspace{0.4cm}
{\Large (sola visualizzazione)\par}
\vspace{1.5cm}
{\large Riepilogo della consegna\par}
\vspace{0.6cm}
{\large Esame ``Progettazione Full Stack''\par}
{\large Prof.\ Bianchini --- UniBS\par}
\vspace{0.4cm}
{\large Appello \textbf{20260629}\par}
\vfill
\begin{minipage}{0.85\textwidth}
\small\itshape
Documento di studio: raccoglie \textbf{tutte} le modifiche fatte rispetto al
progetto originale per soddisfare i tre esercizi della consegna.
\end{minipage}
\vfill
{\large \today\par}
\end{titlepage}

% Tolleranza di impaginazione leggermente più alta: evita stringhe tt che
% sforano nel margine (a scapito di spazi tra parole talvolta più ampi).
\sloppy

\tableofcontents
\newpage

% ====================================================================
\section{La consegna in breve}

Introdurre la \textbf{gestione utenti limitata alla sola visualizzazione}
degli utenti registrati:

\begin{itemize}[leftmargin=1.4em]
  \item Pagina accessibile \textbf{solo agli ADMIN}.
  \item Tabella con colonne \textbf{id, name, email, role}.
  \item Pulsanti \textbf{Nuovo / Modifica / Elimina} (con conferma sulla cancellazione).
  \item Stile uguale alle pagine libri/autori/categorie già esistenti.
  \item Le pagine di \textbf{insert/edit NON sono richieste} (la cancellazione era
        opzionale, ma è stata comunque implementata e funzionante).
\end{itemize}

\paragraph{I tre esercizi.}
\begin{enumerate}[leftmargin=1.6em]
  \item \textbf{Esercizio 1} --- modello ipertestuale \textbf{WebML} della pagina.
  \item \textbf{Esercizio 2} --- rotta \textbf{backend NestJS} che espone gli utenti.
  \item \textbf{Esercizio 3} --- pagina \textbf{frontend React}.
\end{enumerate}

\paragraph{Architettura del progetto.}
Monorepo \textbf{Nx}, con \file{apps/api} (NestJS), \file{apps/ui}
(React + Vite), e librerie \file{libs/server/\{auth,users,books,security\}}.

% ====================================================================
\section{Esercizio 1 --- Modello WebML (documentazione)}

\textbf{File creato:} \file{Appello\_20260629/docs/Esercizio1\_WebML\_GestioneUtenti.md}

Questo esercizio è solo documentale: non tocca il codice, ma \textbf{giustifica}
le scelte di Es.2 ed Es.3 tramite il modello ipertestuale WebML. Di seguito il
contenuto \textbf{integrale} del documento.

\begin{quote}\small\itshape
Progettazione model-driven dei contenuti ipertestuali (WebML).
Riferimenti: slide L9 (notazione) e L10 (running example ``Area libri'' +
mapping all'implementazione).
\end{quote}

% --------------------------------------------------------------------
\subsection{Premesse e requisiti}

Dall'analisi dei requisiti (consegna):

\begin{itemize}[leftmargin=1.4em]
  \item Funzionalità di \textbf{gestione utenti}, limitata alla \textbf{sola
        visualizzazione} degli utenti registrati.
  \item Accesso consentito \textbf{solo agli amministratori} $\rightarrow$ la
        pagina appartiene alla \textbf{site view privata dell'amministratore}
        (area \emph{privata}: accesso tramite autenticazione + autorizzazione di
        ruolo \code{ADMIN}).
  \item La pagina mostra una \textbf{tabella} con \code{id}, \code{nome}
        (\code{name}), \code{email}, \code{ruolo} (\code{role}).
  \item Sono presenti i \textbf{pulsanti} per: inserimento nuovo utente, modifica
        utente, cancellazione (con conferma).
  \item Le pagine di inserimento/modifica \textbf{non} sono realizzate (i relativi
        link puntano a pagine non ancora sviluppate). La cancellazione è opzionale.
\end{itemize}

\textbf{Entità coinvolta} (dal modello dei dati):
\code{User(id, name, email, passwordHash, role)}.
Nella pubblicazione ipertestuale l'attributo \code{passwordHash} \textbf{non}
viene mai esposto.

% --------------------------------------------------------------------
\subsection{Progetto coarse --- Area Utenti}

Si individua una nuova \textbf{area} coesa, dedicata alla gestione degli utenti,
all'interno della site view privata dell'amministratore (analoga alle aree
``libri'', ``autori'', ``categorie'').

\begin{lstlisting}[style=diagram]
:................................................:
: Area Utenti                                    :
:------------------------------------------------:
: Core   (User)                                  :
: Delete (User)            <-- opzionale         :
: [Create (User)]          <-- pagina non realizzata
: [Modify (User)]          <-- pagina non realizzata
:................................................:
\end{lstlisting}

\begin{itemize}[leftmargin=1.4em]
  \item \textbf{Core(User)}: pubblicazione dei contenuti dell'entità core
        \code{User} $\rightarrow$ richiesta dell'esercizio.
  \item \textbf{Delete(User)}: operazione CUD di cancellazione (opzionale).
  \item \textbf{Create(User) / Modify(User)}: previste come funzioni dell'area, ma
        \textbf{non realizzate} in questo esercizio (i pulsanti sono comunque
        mostrati e i link puntano alle relative pagine).
\end{itemize}

\textbf{Visibilità dell'area:} \emph{privata} (accesso previa autenticazione e con
ruolo \code{ADMIN}).

% --------------------------------------------------------------------
\subsection{Progetto dettagliato --- suddivisione in pagine}

L'area si suddivide nelle seguenti pagine (solo la prima è realizzata in questo
esercizio):

\begin{lstlisting}[style=diagram]
:..........................:   :..........................:   :..........................:
: Lista Utenti          [L] :   : Crea Utente           [L] :   : Modifica Utente          :
:--------------------------:   :--------------------------:   :--------------------------:
: Core   (User)            :   : Create (User)            :   : Modify (User)            :
: Delete (User)            :   : (non realizzata)         :   : (non realizzata)         :
:..........................:   :..........................:   :..........................:
\end{lstlisting}

\begin{itemize}[leftmargin=1.4em]
  \item \textbf{Lista Utenti} --- pagina \textbf{Landmark \code{[L]}}: globalmente
        visibile, raggiungibile da ogni altra pagina della site view tramite la
        \emph{navbar} (voce ``Utenti''). È la pagina oggetto dell'esercizio.
  \item \textbf{Crea Utente} / \textbf{Modifica Utente}: pagine non realizzate
        (placeholder di destinazione dei link).
\end{itemize}

% --------------------------------------------------------------------
\subsection{Specifica di dettaglio della pagina «Lista Utenti»}

\subsubsection{Diagramma WebML}

\begin{lstlisting}[style=diagram]
                         PAGINA: Lista Utenti  [L]
   .............................................................................
   :                                                                           :
   :     +-----------------------------+                                       :
   :     |        UtentiUnit           |                                       :
   :     |   [::] MultipleDetails      |----- (modifica) userID ---------------+--->  PAGINA
   :     |                             |   link contestuale                    :     Modifica Utente
   :     |          User               |                                       :     ( User[ID==?] )
   :     |  (id, name, email, role)    |                                       :
   :     +-----------------------------+                                       :
   :            |              \                                               :
   :            |               \  (elimina) userID                            :
   :            |                \  normal link                                :
   :            |                 v                                            :
   :            |          +-------------+      OK (verde)                      :
   :            |          |   Delete    |---------------------> [refresh]      :
   :            |          |   ( - )     |        ricarica la pagina Lista Utenti
   :            |          |   User      |---------------------> ErrorPage      :
   :            |          +-------------+      KO (rosso)                      :
   :            |                                                              :
   :  (nuovo utente) link non contestuale                                      :
   :            +-------------------------------------------------------------+--->  PAGINA
   :                                                                           :     Crea Utente
   :...........................................................................:

   LEGENDA
   [::] MultipleDetails unit   ( - ) Delete operation unit
   --->  link               ----- link contestuale (trasporta userID)
   OK = link verde (successo)   KO = link rosso (fallimento)
\end{lstlisting}

\subsubsection{View component --- \texttt{UtentiUnit} (MultipleDetails)}

Si usa una \textbf{MultipleDetails unit} perché la pagina pubblica \textbf{un
insieme di istanze} dell'entità \code{User}, mostrando per ciascuna più attributi
contemporaneamente (la tabella).

\renewcommand{\arraystretch}{1.3}
\begin{tabularx}{\textwidth}{>{\bfseries}p{4.2cm} X}
\toprule
\normalfont Proprietà & Valore \\
\midrule
Tipo unità & \textbf{MultipleDetails} \\
Contenitore (sorgente) & entità \code{User} \\
Selettore & \emph{nessuno} $\rightarrow$ tutte le istanze (l'admin vede tutti gli utenti) \\
Attributi pubblicati & \code{id}, \code{name}, \code{email}, \code{role} (mai \code{passwordHash}) \\
Parametri in input & nessuno \\
Parametri in output & l'insieme degli \code{\{id\}} degli utenti pubblicati; \code{userID} della riga selezionata per i link \\
\bottomrule
\end{tabularx}

\subsubsection{Operazione --- \texttt{Delete(User)} (opzionale)}

\begin{tabularx}{\textwidth}{>{\bfseries}p{3.6cm} X}
\toprule
\normalfont Proprietà & Valore \\
\midrule
Tipo & \textbf{Delete} operation unit \\
Entità & \code{User} \\
Normal link (in) & dalla \code{UtentiUnit}, trasporta lo \code{userID} dell'utente da eliminare \\
Link OK (verde) & torna a \textbf{Lista Utenti} (ricalcolo/refresh della MultipleDetails) \\
Link KO (rosso) & va a \textbf{ErrorPage} \\
\bottomrule
\end{tabularx}

\medskip
La conferma (``con conferma'') è modellata come interazione che precede
l'attivazione del \emph{normal link} verso l'operazione \code{Delete}.

\subsubsection{Link uscenti dalla pagina}

\begin{tabularx}{\textwidth}{>{\bfseries}p{2.6cm} p{3.6cm} p{3.4cm} X}
\toprule
\normalfont Pulsante & \normalfont Tipo di link & \normalfont Destinazione & \normalfont Contesto trasportato \\
\midrule
Nuovo utente & link \textbf{non contestuale} & pagina Crea Utente & nessuno \\
Modifica & link \textbf{contestuale} & pagina Modifica Utente & \code{userID} ($\rightarrow$ \code{User[ID==?]}) \\
Elimina & \textbf{normal link} verso operazione Delete & Delete(User) & \code{userID} \\
\bottomrule
\end{tabularx}
\renewcommand{\arraystretch}{1}

% --------------------------------------------------------------------
\subsection{Layout di pagina (navbar + utente corrente)}

Come nel running example, il \textbf{Layout page} è un modulo \code{[M]} globale
che definisce navbar e informazioni sull'utente autenticato. La voce di menu
\textbf{``Utenti''} (verso la pagina \emph{Lista Utenti} landmark) è visibile
\textbf{solo} se l'utente corrente ha ruolo \code{ADMIN}.

\begin{lstlisting}[style=diagram]
:..............................:
: Layout page              [M] :
:------------------------------:
:   ( ) CurrentUser            :
:        |                     :
:        v                     :
:   [=] User [ID==?]           :   (Details dell'utente loggato, mostrato in navbar)
:..............................:
   navbar -> { Catalogo | Autori | Categorie | **Utenti (solo ADMIN)** | Logout }
\end{lstlisting}

% --------------------------------------------------------------------
\subsection{Mapping verso l'implementazione (coerente con L10)}

\renewcommand{\arraystretch}{1.3}
\begin{tabularx}{\textwidth}{>{\raggedright\arraybackslash}p{5.2cm} X}
\toprule
\textbf{Costrutto WebML} & \textbf{Implementazione} \\
\midrule
Pagina \emph{Lista Utenti} & componente React \file{users.page.tsx} \\
View component \emph{UtentiUnit} (visualizz.) & tabella renderizzata nel componente React \\
Caricamento dati dal backend & funzione \code{fetchUsers()} in \file{users.api.ts} $\rightarrow$ \code{GET /api/users} (solo ADMIN) \\
Operazione \emph{Delete(User)} & funzione \code{deleteUser(id)} in \file{users.api.ts} $\rightarrow$ \code{DELETE /api/users/:id} \\
Link automatico (rendering) & \code{useEffect} che invoca \code{fetchUsers()} al mount \\
Pagina landmark \emph{Utenti} & voce della \textbf{navbar} nel \file{app-layout.tsx} (visibile solo agli ADMIN) \\
Link contestuale \emph{Modifica} & \code{navigate('/users/:id/edit')} \\
Link non contestuale \emph{Nuovo utente} & \code{navigate('/users/new')} \\
ErrorPage (link KO) & gestione dello stato \code{error} nel componente \\
\bottomrule
\end{tabularx}
\renewcommand{\arraystretch}{1}

% ====================================================================
\section{Esercizio 2 --- Backend NestJS}

\textbf{Obiettivo:} esporre la lista utenti \textbf{senza mai restituire
\code{passwordHash}}, su una rotta protetta accessibile solo agli ADMIN (la rotta
\code{GET /api/users} già esisteva).

% --------------------------------------------------------------------
\subsection{Nuova interfaccia \texttt{UserListItem} (FILE NUOVO)}

\file{libs/server/users/src/lib/interfaces/user-list-item.interface.ts}

\begin{lstlisting}[language=TypeScriptX]
import { UserRole } from '../dto/user-role.enum';

/**
 * Proiezione "sicura" di un utente per la pubblicazione ipertestuale
 * (View component UtentiUnit dell'Esercizio 1): espone solo gli attributi
 * mostrati nella tabella, escludendo sempre `passwordHash`.
 */
export interface UserListItem {
  id: number;
  name: string;
  email: string;
  role: UserRole;
}
\end{lstlisting}

\textbf{Perché:} è la ``proiezione sicura'' dell'entità utente. Definisce un tipo
che contiene \textbf{solo} i campi mostrati in tabella, garantendo a livello di
tipo che \code{passwordHash} non venga propagato verso il frontend.

% --------------------------------------------------------------------
\subsection{Export dell'interfaccia (MODIFICA)}

\file{libs/server/users/src/index.ts} --- aggiunta l'ultima riga:

\begin{lstlisting}[language=TypeScriptX]
export * from './lib/users.module';
export * from './lib/users.service';
export * from './lib/users.repository';
export * from './lib/user.entity';
export * from './lib/dto/user-role.enum';
export * from './lib/dto/create-user.dto';
export * from './lib/dto/update-user.dto';
export * from './lib/interfaces/user-list-item.interface';   // <-- AGGIUNTA
\end{lstlisting}

\textbf{Perché:} rende \code{UserListItem} importabile sia dal backend sia dal
frontend (\code{import \{ UserListItem \} from '@server/users'}), così client e
server condividono lo stesso contratto di tipo.

% --------------------------------------------------------------------
\subsection{Mapping in \texttt{getUsers()} (MODIFICA)}

\file{libs/server/users/src/lib/users.service.ts}

\begin{lstlisting}[language=TypeScriptX]
import { UserListItem } from './interfaces/user-list-item.interface';   // <-- import aggiunto

async getUsers(role?: UserRole): Promise<UserListItem[]> {              // <-- tipo di ritorno cambiato
    const users = await this.usersRepository.findAll(role);

    if (role && users.length === 0) {
        throw new NotFoundException(`No users found with role ${role}`);
    }

    // Mapping verso la proiezione "sicura": non esporre mai passwordHash.
    return users.map(({ id, name, email, role }) => ({ id, name, email, role }));   // <-- AGGIUNTA
}
\end{lstlisting}

\textbf{Perché:} prima il metodo restituiva le entità complete (incluso
\code{passwordHash}). Ora il \code{.map()} estrae solo i quattro campi previsti
$\rightarrow$ la rotta non espone più dati sensibili.

% --------------------------------------------------------------------
\subsection{Quadro completo delle rotte del controller utenti}

\file{libs/server/users/src/lib/users.controller.ts} --- il controller è montato su
\code{@Controller('users')} e il prefisso globale dell'app è \code{/api}, quindi
tutte le rotte vivono sotto \code{/api/users}. Ecco \textbf{ogni rotta} con le
relative protezioni:

\renewcommand{\arraystretch}{1.3}
\small
\begin{longtable}{>{\ttfamily}p{2.9cm} p{2.4cm} p{1.9cm} p{3.2cm} p{2.2cm}}
\toprule
\normalfont\textbf{Metodo + path} & \textbf{Guardie} & \textbf{Ruoli ammessi} & \textbf{Cosa fa} & \textbf{Usata?} \\
\midrule
\endhead
GET /users (opz. ?role=) & JwtAuthGuard, RolesGuard & \textbf{ADMIN} & Lista utenti (proiezione sicura \code{UserListItem[]}) & \si{} --- alimenta la tabella \\
GET /users/me & JwtAuthGuard & qualsiasi loggato & Restituisce l'utente corrente dal token & indirettamente (navbar) \\
GET /users/interns & nessuna & pubblica & Placeholder (``API non implementata'') & no \\
GET /users/:id & JwtAuthGuard, RolesGuard & ADMIN, USER & Dettaglio singolo utente & no \\
POST /users & nessuna & pubblica & Crea utente (accetta \code{role}) & no (preesistente) \\
PATCH /users/:id & JwtAuthGuard, RolesGuard & ADMIN & Aggiorna utente & no \\
DELETE /users/:id & JwtAuthGuard, RolesGuard & ADMIN & Cancella utente & \si{} --- pulsante ``Elimina'' \\
\bottomrule
\end{longtable}
\normalsize
\renewcommand{\arraystretch}{1}

\begin{quote}\small
\textbf{Importante:} l'ordine di dichiarazione conta. \code{GET /users/me} e
\code{GET /users/interns} sono dichiarati \textbf{prima} di \code{GET /users/:id},
altrimenti \code{:id} catturerebbe anche \code{me}/\code{interns} interpretandoli
come id.
\end{quote}

\textbf{Le due rotte effettivamente sfruttate dalla pagina:}

\begin{lstlisting}[language=TypeScriptX]
// 1) Lettura della lista (tabella)
@Get() // GET /api/users  oppure  /api/users?role=value
@UseGuards(JwtAuthGuard, RolesGuard)
@Roles(UserRole.ADMIN)
@ApiBearerAuth()
@ApiQuery({ name: 'role', required: false, enum: UserRole })
getUsers(@Query('role', new ParseEnumPipe(UserRole, { optional: true })) role?: UserRole) {
    return this.serverUsersService.getUsers(role);
}

// 2) Cancellazione (pulsante Elimina)
@Delete(':id') // DELETE /api/users/:id
@UseGuards(JwtAuthGuard, RolesGuard)
@Roles(UserRole.ADMIN)
@ApiBearerAuth()
removeUser(@Param('id', ParseIntPipe) id: number) {
    return this.serverUsersService.removeUser(id);
}
\end{lstlisting}

\textbf{Come funziona la protezione (catena di guardie):}

\begin{enumerate}[leftmargin=1.6em]
  \item \code{JwtAuthGuard} --- legge l'header \code{Authorization: Bearer <token>},
        valida il JWT e popola \code{request.user}. Senza token valido
        $\rightarrow$ \textbf{401 Unauthorized}.
  \item \code{RolesGuard} --- legge i ruoli richiesti impostati dal decoratore
        \code{@Roles(...)} (via metadata Reflector) e li confronta con il ruolo
        dell'utente nel token. Se non combaciano $\rightarrow$ \textbf{403 Forbidden}.
  \item \code{@Roles(UserRole.ADMIN)} --- non esegue logica: \textbf{annota} la rotta
        con i ruoli ammessi, che \code{RolesGuard} poi legge. È il pezzo che impone
        ``solo ADMIN''.
  \item \code{@ApiBearerAuth()} / \code{@ApiQuery(...)} --- solo documentazione
        Swagger, nessun effetto a runtime.
\end{enumerate}

\textbf{Pipe sui parametri:}

\begin{itemize}[leftmargin=1.4em]
  \item \code{ParseIntPipe} su \code{:id} $\rightarrow$ converte e valida che l'id
        sia un intero (altrimenti 400).
  \item \code{ParseEnumPipe(UserRole, \{ optional: true \})} sul query \code{role}
        $\rightarrow$ accetta solo valori validi dell'enum, ed è opzionale.
\end{itemize}

\textbf{Flusso completo di una richiesta \code{GET /api/users} (end-to-end):}

\begin{lstlisting}[style=diagram]
UsersPage (useEffect)
   -> fetchUsers()                       [users.api.ts]
   -> GET /api/users  + Bearer token
        -> JwtAuthGuard  (valida token, set request.user)
        -> RolesGuard    (richiede ADMIN)
        -> ServerUsersController.getUsers()
        -> ServerUsersService.getUsers()  [mappa a UserListItem, niente passwordHash]
        -> UsersRepository.findAll()
   <- JSON: UserListItem[]
   -> setUsers(...)  -> render tabella
\end{lstlisting}

% ====================================================================
\section{Esercizio 3 --- Frontend React}

\textbf{Obiettivo:} nuova pagina \code{/users} (solo ADMIN) con tabella e pulsanti,
in stile identico alle altre pagine.

% --------------------------------------------------------------------
\subsection{API client (FILE NUOVO)}

\file{apps/ui/src/features/users/users.api.ts}

\textbf{Codice completo del file (fedele all'originale):}

\begin{lstlisting}[language=TypeScriptX]
import { UserListItem } from '@server/users';
import { handleApiError } from '../shared/utils.api';

const API_URL = 'http://localhost:3333/api';

function getAuthHeaders() {
    const token = localStorage.getItem('access_token');

    return {
        'Content-Type': 'application/json',
        Authorization: `Bearer ${token}`,
    };
}

export async function fetchUsers(): Promise<UserListItem[]> {
    const response = await fetch(`${API_URL}/users`, {
        headers: getAuthHeaders(),
    });

    if (!response.ok) {
        await handleApiError(response);
    }

    return response.json();
}

export async function deleteUser(id: number): Promise<void> {
    const response = await fetch(`${API_URL}/users/${id}`, {
        method: 'DELETE',
        headers: getAuthHeaders(),
    });

    if (!response.ok) {
        await handleApiError(response);
    }
}
\end{lstlisting}

\textbf{Perché:} centralizza le chiamate HTTP (con token JWT negli header) e riusa
il tipo \code{UserListItem} esportato dal backend. Coerente con gli altri
\file{*.api.ts}.

% --------------------------------------------------------------------
\subsection{Pagina elenco utenti (FILE NUOVO)}

\file{apps/ui/src/features/users/users.page.tsx}

\textbf{Punti chiave:}
\begin{itemize}[leftmargin=1.4em]
  \item Carica gli utenti con \code{fetchUsers()} dentro un \code{useEffect},
        gestendo gli stati \code{loading} / \code{error}.
  \item Tabella con colonne \textbf{Id / Nome / Email / Ruolo / Azioni}.
  \item Pulsante \textbf{``Nuovo utente''} $\rightarrow$ naviga a \code{/users/new}.
  \item Pulsante \textbf{``Modifica''} per riga $\rightarrow$ naviga a
        \code{/users/:id/edit}.
  \item Pulsante \textbf{``Elimina''} per riga $\rightarrow$ \code{window.confirm}
        di conferma, poi \code{deleteUser(id)} e aggiornamento ottimistico della
        lista (\code{filter}).
  \item Stile riutilizzato da \file{../css/books.module.css}.
\end{itemize}

\textbf{Codice completo del file:}

\begin{lstlisting}[language=TypeScriptX]
import { useEffect, useState } from 'react';
import { useNavigate } from 'react-router-dom';
import clsx from 'clsx';
import book_styles from '../css/books.module.css';
import { UserListItem } from '@server/users';
import { deleteUser, fetchUsers } from './users.api';

export function UsersPage() {
  const navigate = useNavigate();
  const [users, setUsers] = useState<UserListItem[]>([]);
  const [loading, setLoading] = useState(true);
  const [error, setError] = useState<string | null>(null);

  useEffect(() => {
    fetchUsers()
      .then(setUsers)
      .catch((err) => setError(err.message))
      .finally(() => setLoading(false));
  }, []);

  async function handleDelete(id: number) {
    const confirmed = window.confirm('Vuoi davvero cancellare questo utente?');

    if (!confirmed) return;

    try {
      await deleteUser(id);
      setUsers((current) => current.filter((user) => user.id !== id));
    } catch (err: any) {
      setError(err.message);
    }
  }

  if (loading) {
    return (
      <main className={book_styles.page}>
        <section className={book_styles.card}>
          <p className={book_styles.message}>Caricamento utenti...</p>
        </section>
      </main>
    );
  }

  if (error) {
    return (
      <main className={book_styles.page}>
        <section className={book_styles.card}>
          <p className={book_styles.error}>{error}</p>
        </section>
      </main>
    );
  }

  return (
    <main className={book_styles.page}>
      <section className={clsx(book_styles.card, book_styles.cardLarge)}>
        <header className={book_styles.header}>
          <h1 className={book_styles.title}>👤 Utenti</h1>
          <p className={book_styles.subtitle}>
            Elenco degli utenti registrati nella libreria.
          </p>
        </header>

        <button
          className={book_styles.button}
          onClick={() => navigate('/users/new')}
        >
          Nuovo utente
        </button>

        {users.length === 0 ? (
          <p className={book_styles.message}>Nessun utente disponibile.</p>
        ) : (
          <div className={book_styles.tableWrapper}>
            <table className={book_styles.table}>
              <thead>
                <tr>
                  <th className={book_styles.th}>Id</th>
                  <th className={book_styles.th}>Nome</th>
                  <th className={book_styles.th}>Email</th>
                  <th className={book_styles.th}>Ruolo</th>
                  <th className={book_styles.th}>Azioni</th>
                </tr>
              </thead>

              <tbody>
                {users.map((user) => (
                  <tr key={user.id} className={book_styles.row}>
                    <td className={book_styles.td}>{user.id}</td>
                    <td className={book_styles.titleCell}>{user.name}</td>
                    <td className={book_styles.td}>{user.email}</td>
                    <td className={book_styles.td}>{user.role}</td>
                    <td className={book_styles.td}>
                      <button
                        className={book_styles.secondaryButton}
                        onClick={() => navigate(`/users/${user.id}/edit`)}
                      >
                        Modifica
                      </button>

                      <button
                        className={book_styles.dangerButton}
                        onClick={() => handleDelete(user.id)}
                      >
                        Elimina
                      </button>
                    </td>
                  </tr>
                ))}
              </tbody>
            </table>
          </div>
        )}
      </section>
    </main>
  );
}
\end{lstlisting}

\textbf{Spiegazione riga per riga delle scelte:}

\begin{itemize}[leftmargin=1.4em]
  \item \code{const [users, setUsers] = useState<UserListItem[]>([])} --- lo stato è
        \textbf{tipizzato} con il contratto condiviso \code{UserListItem}: il
        compilatore garantisce che si accede solo a \code{id/name/email/role}.
  \item \code{useEffect(..., [])} con array di dipendenze vuoto $\rightarrow$ il
        fetch parte \textbf{una sola volta} al montaggio del componente.
  \item \code{.then(setSomething).catch(...).finally(...)} $\rightarrow$ pattern a
        tre stati (dati / errore / caricamento), identico alle altre pagine del
        progetto.
  \item \code{handleDelete} $\rightarrow$ \textbf{conferma} con \code{window.confirm}
        (requisito della consegna), poi cancella sul server e infine aggiorna lo
        stato locale con \code{filter((user) => user.id !== id)}
        (\textbf{aggiornamento ottimistico}: non si rifà la fetch, si rimuove la
        riga dallo stato).
  \item Rendering condizionale: \code{loading} $\rightarrow$ messaggio; \code{error}
        $\rightarrow$ messaggio d'errore; lista vuota $\rightarrow$ ``Nessun utente
        disponibile''; altrimenti $\rightarrow$ tabella.
  \item Le classi CSS (\code{page}, \code{card}, \code{cardLarge}, \code{table},
        \code{th}, \code{td}, \code{titleCell}, \code{secondaryButton},
        \code{dangerButton}, \ldots) sono \textbf{riusate} da
        \file{books.module.css}, così lo stile è identico alle pagine
        libri/autori/categorie.
\end{itemize}

\begin{quote}\small
\textbf{Nota sull'evoluzione:} nella prima versione i pulsanti ``Nuovo'' e
``Modifica'' erano \code{disabled} con \code{title="Funzione non disponibile"}. In
un secondo momento sono stati resi cliccabili e fatti puntare alla paginetta
``Funzione non implementata'' (vedi sezione successiva).
\end{quote}

% --------------------------------------------------------------------
\subsection{Pagina ``Funzione non implementata'' (FILE NUOVO)}

\file{apps/ui/src/features/users/user-not-implemented.page.tsx}

\begin{lstlisting}[language=TypeScriptX]
import { useNavigate } from 'react-router-dom';
import clsx from 'clsx';
import book_styles from '../css/books.module.css';

export function UserNotImplementedPage() {
  const navigate = useNavigate();
  return (
    <main className={book_styles.page}>
      <section className={clsx(book_styles.card, book_styles.cardSmall)}>
        <h1 className={book_styles.title}>🚧 Funzione non implementata</h1>
        <p className={book_styles.message}>
          La creazione e la modifica degli utenti non sono ancora disponibili.
          La gestione utenti è attualmente limitata alla sola visualizzazione.
        </p>
        <button className={book_styles.button} onClick={() => navigate('/users')}>
          Torna agli utenti
        </button>
      </section>
    </main>
  );
}
\end{lstlisting}

\textbf{Perché:} la consegna non richiede insert/edit. Invece di lasciare pulsanti
morti, i pulsanti Nuovo/Modifica portano a questa pagina che dichiara
esplicitamente che la funzione non è implementata, mantenendo la UX coerente.

% --------------------------------------------------------------------
\subsection{Rotte (MODIFICA)}

\file{apps/ui/src/app/app.tsx} --- sono stati aggiunti \textbf{2 import} e
\textbf{3 rotte} (le righe marcate \code{// <-- AGGIUNTA}). Di seguito il
\textbf{file completo} per vedere le rotte nel loro contesto:

\begin{lstlisting}[language=TypeScriptX]
import { BooksPage } from '../features/books/books.page';
import { Routes, Route, Navigate } from 'react-router-dom';
import { LoginPage } from '../features/auth/login.page';
import { ProtectedRoute } from '../features/auth/protected-route';
import { LogoutPage } from '../features/auth/logout.page';
import { CreateBookPage } from '../features/books/create-book.page';
import { EditBookPage } from '../features/books/edit-book.page';
import { AppLayout } from '../features/layouts/app-layout';
import { HomeModulePage } from '../features/books/home-module.page';
import { HomeTailwindPage } from '../features/books/home-tailwind.page';
import { HomeBootstrapPage } from '../features/books/home-bootstrap.page';
import { AuthorsPage } from '../features/authors/authors.page';
import { CreateAuthorPage } from '../features/authors/create-author.page';
import { EditAuthorPage } from '../features/authors/edit-author.page';
import { CategoriesPage } from '../features/categories/categories.page';
import { CreateCategoryPage } from '../features/categories/create-category.page';
import { EditCategoryPage } from '../features/categories/edit-category.page';
import { UsersPage } from '../features/users/users.page';                              // <-- AGGIUNTA
import { UserNotImplementedPage } from '../features/users/user-not-implemented.page';  // <-- AGGIUNTA

export function App() {
  return (
    <Routes>
      <Route path="/" element={<Navigate to="/books" replace />} />
      <Route path="/login" element={<LoginPage />} />
      <Route path="/logout" element={<LogoutPage />} />

      <Route
        element={
          <ProtectedRoute>
            <AppLayout />
          </ProtectedRoute>
        }
      >
        <Route path="/books" element={<BooksPage />} />
        <Route path="/books/new" element={<CreateBookPage />} />
        <Route path="/books/:id/edit" element={<EditBookPage />} />
        <Route path="/authors" element={<AuthorsPage />} />
        <Route path="/authors/new" element={<CreateAuthorPage />} />
        <Route path="/authors/:id/edit" element={<EditAuthorPage />} />
        <Route path="/categories" element={<CategoriesPage />} />
        <Route path="/categories/new" element={<CreateCategoryPage />} />
        <Route path="/categories/:id/edit" element={<EditCategoryPage />} />
        <Route path="/users" element={<UsersPage />} />                       {/* <-- AGGIUNTA */}
        <Route path="/users/new" element={<UserNotImplementedPage />} />      {/* <-- AGGIUNTA */}
        <Route path="/users/:id/edit" element={<UserNotImplementedPage />} /> {/* <-- AGGIUNTA */}
      </Route>

      <Route path="/home-module" element={<HomeModulePage />} />
      <Route path="/home-tailwind" element={<HomeTailwindPage />} />
      <Route path="/home-bootstrap" element={<HomeBootstrapPage />} />
    </Routes>
  );
}

export default App;
\end{lstlisting}

\textbf{Perché:} le tre rotte utenti sono inserite \textbf{dentro} il blocco
protetto da \code{<ProtectedRoute>\allowbreak<AppLayout/>\allowbreak</ProtectedRoute>} (richiedono login
e mostrano la navbar). \code{/users/new} e \code{/users/:id/edit} riusano la stessa
pagina ``non implementata''. Le rotte pubbliche (\code{/login}, \code{/logout}, le
home demo) restano fuori dal blocco protetto.

% --------------------------------------------------------------------
\subsection{Voce di menu nella navbar (MODIFICA)}

\file{apps/ui/src/features/layouts/app-layout.tsx} --- aggiunta la voce ``Utenti''
nella sezione \code{navLinks} (riga marcata). Contesto del blocco modificato:

\begin{lstlisting}[language=TypeScriptX]
<div className={styles.navLinks}>
  <button onClick={() => navigate('/books')}>Catalogo</button>
  <button onClick={() => navigate('/books/new')}>Nuovo libro</button>
  <button onClick={() => navigate('/authors')}>Autori</button>
  <button onClick={() => navigate('/categories')}>Categorie</button>
  {user?.role === 'ADMIN' && (                                      {/* <-- AGGIUNTA */}
    <button onClick={() => navigate('/users')}>Utenti</button>      {/* <-- AGGIUNTA */}
  )}                                                                {/* <-- AGGIUNTA */}

  <div className={styles.userSection}>
    {/* ... pulsante utente + dropdown logout (invariati) ... */}
  </div>
</div>
\end{lstlisting}

Il ruolo arriva dallo stato \code{user}, popolato al mount da
\code{fetchCurrentUser()} (\code{GET /api/users/me}):

\begin{lstlisting}[language=TypeScriptX]
const [user, setUser] = useState(null);

useEffect(() => {
  fetchCurrentUser()
    .then(setUser)
    .catch((err) => setUser(null));
}, []);
\end{lstlisting}

\textbf{Perché:} la voce ``Utenti'' compare nel menu \textbf{solo se l'utente
loggato è ADMIN} (\code{user?.role === 'ADMIN'}), in linea con la consegna (pagina
riservata agli ADMIN).

% --------------------------------------------------------------------
\subsection{Approfondimento --- come funziona il routing frontend}

Le tre rotte utenti sono \textbf{annidate} dentro una rotta ``wrapper'' senza
\code{path}:

\begin{lstlisting}[language=TypeScriptX]
<Route
  element={
    <ProtectedRoute>
      <AppLayout />
    </ProtectedRoute>
  }
>
  {/* ... books, authors, categories ... */}
  <Route path="/users" element={<UsersPage />} />
  <Route path="/users/new" element={<UserNotImplementedPage />} />
  <Route path="/users/:id/edit" element={<UserNotImplementedPage />} />
</Route>
\end{lstlisting}

Cosa implica, passo per passo:

\begin{enumerate}[leftmargin=1.6em]
  \item \textbf{Rotta wrapper (layout route):} non ha \code{path}, serve solo a
        raggruppare le rotte figlie e ad applicare loro lo stesso ``guscio''.
  \item \code{<ProtectedRoute>} --- verifica l'autenticazione (presenza/validità del
        token); se l'utente non è loggato reindirizza al login. Tutte le rotte
        figlie ereditano questa protezione $\rightarrow$ \code{/users},
        \code{/users/new}, \code{/users/:id/edit} sono \textbf{tutte} accessibili
        solo da utenti autenticati.
  \item \code{<AppLayout>} --- disegna la navbar e poi un \code{<Outlet/>}: è dentro
        l'\code{Outlet} che React Router inietta la pagina figlia corrente. Per
        questo la navbar resta fissa mentre cambia solo il contenuto.
  \item \code{navigate(...)} --- la navigazione è programmatica (hook
        \code{useNavigate}):
        \begin{itemize}[leftmargin=1.2em]
          \item \code{navigate('/users')} dalla navbar (solo ADMIN);
          \item \code{navigate('/users/new')} dal pulsante ``Nuovo utente'';
          \item \code{navigate('/users/\$\{user.id\}/edit')} dal pulsante
                ``Modifica'' (URL dinamico con id);
          \item \code{navigate('/users')} dal pulsante ``Torna agli utenti'' nella
                pagina non implementata.
        \end{itemize}
  \item \code{:id} è un parametro dinamico: qualsiasi valore (\code{/users/7/edit},
        \code{/users/42/edit}, \ldots) matcha la stessa rotta. Qui punta alla pagina
        ``non implementata'', quindi l'id non viene letto; in una vera edit si
        leggerebbe con \code{useParams()}.
\end{enumerate}

\textbf{Livelli di protezione per \code{/users} (tre livelli):}

\renewcommand{\arraystretch}{1.3}
\begin{tabularx}{\textwidth}{>{\bfseries}p{2.6cm} p{4.6cm} X}
\toprule
\normalfont Livello & \normalfont Meccanismo & \normalfont Cosa blocca \\
\midrule
Navbar & \code{user?.role === 'ADMIN'} & nasconde il link ai non-ADMIN \\
Frontend route & \code{<ProtectedRoute>} & blocca gli utenti non autenticati \\
Backend & \code{JwtAuthGuard} + \code{RolesGuard} + \code{@Roles(ADMIN)} & 401/403 su accesso non autorizzato (anche via URL diretto) \\
\bottomrule
\end{tabularx}
\renewcommand{\arraystretch}{1}

\medskip
\textbf{Nota:} la rotta frontend \code{/users} controlla solo l'autenticazione, non
il ruolo; la vera barriera sul ruolo è il \textbf{backend}, che risponde 403 a un
non-ADMIN che forzasse l'URL.

% ====================================================================
\section{Configurazione Nx}

Avendo introdotto una dipendenza nuova \code{ui $\rightarrow$ @server/users} (per
importare il tipo \code{UserListItem}), è stato eseguito \code{nx sync} per
aggiornare i \emph{project references} di TypeScript.

% ====================================================================
\section{Verifica finale}

\begin{itemize}[leftmargin=1.4em]
  \item \code{nx run-many -t build -p api ui} $\rightarrow$ \textbf{build OK} per
        entrambi i progetti.
  \item \code{nx lint}:
        \begin{itemize}[leftmargin=1.2em]
          \item \code{ui} $\rightarrow$ solo \emph{warning} preesistenti (\code{any}
                + emoji), identici a \code{categories.page}.
          \item \code{@server/\allowbreak users:lint} $\rightarrow$ fallisce per una
                \textbf{dipendenza circolare PREESISTENTE} \code{@server/users}
                $\leftrightarrow$ \code{@server/security} (controller $\leftrightarrow$
                roles.guard). \textbf{Non introdotta da noi.}
        \end{itemize}
\end{itemize}

% ====================================================================
\section{Note / limiti noti (fuori scope della consegna)}

\begin{itemize}[leftmargin=1.4em]
  \item La pagina \code{/users} è protetta da \code{ProtectedRoute}
        (autenticazione) ma \textbf{non} ricontrolla il ruolo a livello di rotta
        frontend: un non-ADMIN che digita \code{/users} a mano riceve comunque
        \textbf{403 dal backend} (la pagina mostra l'errore). La voce in navbar è
        già nascosta ai non-ADMIN.
  \item \code{POST /users} e \code{register} restano \textbf{pubbliche} e accettano
        \code{role} (possibile \emph{privilege escalation}) --- comportamento
        preesistente, fuori dallo scope.
  \item Dipendenza circolare tra librerie server --- preesistente.
\end{itemize}

% ====================================================================
\section{Elenco rapido dei file toccati}

\renewcommand{\arraystretch}{1.3}
\begin{tabularx}{\textwidth}{X p{2.4cm} p{1.6cm}}
\toprule
\textbf{File} & \textbf{Tipo} & \textbf{Esercizio} \\
\midrule
\file{docs/Esercizio1\_WebML\_GestioneUtenti.md} & nuovo & Es.1 \\
\file{libs/server/users/src/lib/interfaces/user-list-item.interface.ts} & nuovo & Es.2 \\
\file{libs/server/users/src/index.ts} & modificato & Es.2 \\
\file{libs/server/users/src/lib/users.service.ts} & modificato & Es.2 \\
\file{apps/ui/src/features/users/users.api.ts} & nuovo & Es.3 \\
\file{apps/ui/src/features/users/users.page.tsx} & nuovo & Es.3 \\
\file{apps/ui/src/features/users/user-not-implemented.page.tsx} & nuovo & Es.3 \\
\file{apps/ui/src/app/app.tsx} & modificato & Es.3 \\
\file{apps/ui/src/features/layouts/app-layout.tsx} & modificato & Es.3 \\
\bottomrule
\end{tabularx}
\renewcommand{\arraystretch}{1}

\end{document}
